%% =====================================================================
%%  B7-T-23-05 -- what B7 contributes to "The Five-Second Window"
%%  -------------------------------------------------------------------
%%  LAYER A: APPEND-ONLY. Written once at the entry's write, then left
%%  alone. If the source has more to say later, add a bullet -- do not
%%  rewrite. Material found ONLY here goes in the `unique` slot with
%%  the unique flag set, which makes it undeletable (rule R10).
%% =====================================================================
%% @kind: source
%% @id: B7-T-23-05
%% @book: B7
%% @topic: T-23-05
%% @locator: chs. 2--3, 8, 12 (pp. 17--52, 57--62, 85--96)
%% @status: distilled
%% @unique: yes
%% @integrated: yes
%% @updated: 2026-09-19

\dossier{B7}{T-23-05}{\bookshort{B7}}{chs. 2--3, 8, 12 (pp. 17--52, 57--62, 85--96)}

\begin{distilled}
  \item No such thing as a human lie detector: nothing in behaviour proves a lie unless it contradicts what you independently know; the methodology raises confidence, never to certainty.
  \item Global behaviour assessment (reading overall manner) fails because the reader must guess the cause of a behaviour; nerves, habit and baseline manner all mimic signals.
  \item The fix is timing: treat the question as the stimulus and look for the first deceptive sign inside the following five seconds; beyond that span the honest mind has moved on (speech ~125--150 words/minute, thought an order faster), so late behaviours belong to something else.
  \item Cluster rule: two or more indicators, verbal or nonverbal in any mix, make a reading; a single indicator is ignored --- habits exist, and confidence rises with the count.
  \item Dual-channel capture (L-squared mode): deliberately look and listen at once; the brain resists it and defaults to one channel unless trained.
  \item Nonverbal indicator catalogue (ch. 8): behavioural pause or delay; verbal/nonverbal disconnect; hiding mouth or eyes; throat-clearing or a significant swallow before answering; hand-to-face activity; anchor-point movement; grooming gestures.
  \item Truthful responses have their own signature (ch. 9): a reader watching only for the deceptive list will convert ordinary nerves into guilt.
  \item Limits (ch. 12): single-sign reading (the television microexpression school) produces false positives at scale; microexpressions are information about emotion, not verdicts about lies.
\end{distilled}
\begin{unique}
  \item The five-second stimulus rule, the cluster rule with its single-indicator exemption, and the L-squared dual-capture discipline are unique to B7 in this corpus.
  \item The nonverbal indicator catalogue (behavioural pause, mouth/eye hiding, throat-clear, hand-to-face, anchor-point movement, grooming) is unique to B7.
\end{unique}
